% Conclusiones del informe del proyecto final (modelo original BDL).
\section{Conclusiones}\label{sec:conclusiones}

El análisis exploratorio mostró que el puntaje de COFINFAD es un índice de pocas variables más un escalón constante que
depende de una variable que hay que excluir. BDL convierte ese hallazgo en un modelo: una suma de árboles más una
mezcla latente de desplazamientos, ajustada con una regla de aprendizaje propia, un EM generalizado cuyo paso M es un
paso de \emph{boosting} contra el objetivo sin escalón, programada desde cero y con la garantía de que la verosimilitud
no decrece.

En COFINFAD, BDL es el modelo más preciso: el menor error en las cuatro métricas, tanto en validación cruzada como en
prueba, y a $0{,}00001$ de RMSE del predictor óptimo. Su ventaja sobre los mejores \emph{boostings} es estadísticamente
detectable en validación cruzada, pero pequeña en la práctica, y no alcanza las tres pruebas que exige el protocolo
para declararlo superior. Tampoco se generaliza: en los cuatro conjuntos adicionales no hay un escalón nítido y BDL no
supera al mejor benchmark de ninguno. En clasificación queda al nivel de los mejores, sin diferencias significativas.

Su valor está en otra parte. Recupera la distribución de la variable excluida sin verla, con cinco clases cuyos pesos
coinciden con las proporciones reales; necesita diez veces menos datos para alcanzar el error de XGBoost, y varía
menos entre semillas que XGBoost, LightGBM, la EBM y Random Forest. La ablación confirma que la mezcla aporta en validación cruzada
y que su aporte crece cuando hay pocos datos.

Las principales limitaciones son que los datos son sintéticos, que cinco conjuntos dan poca potencia a las
comparaciones entre conjuntos y que algunas pruebas se añadieron después del protocolo. Como trabajo futuro quedan
hacer que los pesos de las clases dependan de los predictores, permitir una dispersión distinta por clase y validar el
modelo con datos reales y con fuga observada.
